\chapter{Repository and Reproduction Instructions}
\label{app:reproducibility}
\begingroup\raggedright

The code, configurations and run records are in the project repository,
\url{https://github.com/Jschechter18/fall-2026-group3}. Table~\ref{tab:pipeline-stages}
gives an overview of the pipeline stages. This appendix gives the commands for
each stage, the folder of the final run each one produced, and what has to be
retrieved separately because it is not in Git. The software environment,
seeds and provenance records are described in the last section of the
methodology chapter, and all settings are listed in
Appendix~\ref{app:hyperparameters}.

\section{Setup}

{\footnotesize
\begin{verbatim}
git clone https://github.com/Jschechter18/fall-2026-group3.git
cd fall-2026-group3
conda env create -f environment.yml   # Python 3.11, environment "capstone"
conda activate capstone
python -m pip install -e .
export PYTHONPATH=src
\end{verbatim}
}

The models are read from Hugging Face at the pinned revisions. Synthetic
generation and every labeling step call Claude through Amazon Bedrock.
Collection and the intervention runs need a GPU with at least 20~GB of memory.
Setup details are in \path{cookbooks/setup.md}.

\section{What the Repository Contains}

The repository holds the code (\path{src/mas_sae/} and \path{scripts/}), the
configurations (\path{configs/}), the run records and smaller results
(\path{results/}: labels, probe and feature-selection outputs, intervention
tables and labeled answers), the analysis notebooks (\path{notebooks/}) and
this report (\path{docs/final_report/}). Large files are not in Git: the
MuSiQue copies and synthetic records (\path{data/}), the collected episodes
and activations (\path{artifacts/collection/mixed_final_001/}), SAE weights
and activation caches. They are kept in the project's artifact storage and
must be retrieved before the stages that read them. Every run record stores
the hashes of its inputs, so a retrieved file can be checked against the run
that used it.

\section{Reproducing the Study}

Each stage reads the outputs of the previous one. Partitions are fixed in
stage~1 and respected by every later stage.

\paragraph{0. Data.} MuSiQue is downloaded once. Synthetic generation makes
paid API calls and is run once.

{\footnotesize
\begin{verbatim}
python scripts/data/get_musique_dataset.py
python -u scripts/data/generate_synthetic_musique.py --config configs/data/opus5_v3_full.yaml
\end{verbatim}
}
\noindent Final output: \path{data/MuSiQue/synthetic/runs/opus5-v3-full-001/}.

\paragraph{1. Collection.} \texttt{<run\_dir>/collection.yaml} is a copy of
\path{configs/collection/distributed_final.yaml}. Each worker collects its own
chunks.

{\footnotesize
\begin{verbatim}
python scripts/collection/distributed_collection.py prepare \
    --config <run_dir>/collection.yaml --num-chunks <N> --output <run_dir>/manifest.json
python scripts/collection/collect_activations.py \
    --config <run_dir>/collection.yaml --manifest <run_dir>/manifest.json --chunk <K ...>
python scripts/collection/distributed_collection.py merge --manifest <run_dir>/manifest.json
python scripts/collection/distributed_collection.py split \
    --manifest <run_dir>/manifest.json --config <run_dir>/collection.yaml
\end{verbatim}
}
\noindent Final output: \path{artifacts/collection/mixed_final_001/prepared/}.

\paragraph{2. Behavior labels.} One labeling run per partition, then an export
of the train and validation labels for the probes.

{\footnotesize
\begin{verbatim}
R=artifacts/collection/mixed_final_001/prepared/results/mixed_final_001
L=results/claude_validation/mixed_final_001
python scripts/labels/run_claude_validation.py \
    --episodes $R/<partition>/interactions.jsonl \
    --partition <partition> --output $L/<partition>
python scripts/labels/export_final_labels_for_probes.py \
    --claude-runs $L/train $L/validation --full-results $R \
    --layout partitioned --output results/behavior/mixed_final_001_claude_probe
\end{verbatim}
}
\noindent Final outputs: \path{results/claude_validation/mixed_final_001/}
(one folder per partition) and \path{results/behavior/mixed_final_001_claude_probe/}.

\paragraph{3. Layer selection.}

{\footnotesize
\begin{verbatim}
python scripts/layer_selection/layer_selection.py \
    --config configs/layer_selection/layer_selection_mixed_final_001.yaml
\end{verbatim}
}
\noindent Final output:
\path{results/layer_selection_mixed_final_001_quick/20261102T201434082931Z/}.
This run used a copy of this configuration with the label-shuffle control
switched off (\texttt{n\_permutations: 0}); the copy is not in the repository.

\paragraph{4. SAEs.} The custom SAEs are tuned on the baseline grid, and the
chosen trial is trained to completion. The pretrained SAE is downloaded and
then adapted; adding \texttt{--train full} gives the exploratory full
fine-tune.

{\footnotesize
\begin{verbatim}
python scripts/sae/train_sae.py --hp-tune \
    --config configs/sae/layer_25_mixed_final_001_baselines.yaml
python scripts/sae/train_sae.py --epochs 200 \
    --finish-run results/sae/musique/tuning/layer_25/<study_id>/trials/<trial>
P=results/sae/pretrained/gemma-scope-2-4b-it
A=artifacts/collection/mixed_final_001/prepared/activations/mixed_final_001
hf download google/gemma-scope-2-4b-it \
    --include "resid_post_all/layer_25_width_16k_l0_big/*" --local-dir $P
python scripts/sae/finetune_pretrained_sae.py \
    --sae-dir $P/resid_post_all/layer_25_width_16k_l0_big \
    --activations-dir $A/layer_25 \
    --layer 25 \
    --probe-config configs/feature_selection/feature_selection_mixed_final_001_layer25.yaml
\end{verbatim}
}
\noindent Final outputs:
\begin{itemize}
    \item $L_1$ SAE: \path{results/sae/musique/runs/layer_25/20261102T182346930753Z-150e78c341c6}
    \item TopK SAE: \path{results/sae/musique/runs/layer_25/20261102T182749413400Z-150e78c341c6}
    \item Decoder-adapted SAE: \path{results/sae/decoder_finetuning/20261104T164310458065Z/best}
    \item Fully fine-tuned SAE: \path{results/sae/full_finetuning/20261104T151230848807Z/best}
\end{itemize}

\paragraph{5. Feature selection.} The second command adds the full pool of
random draws to a copy of the run's feature sets.

{\footnotesize
\begin{verbatim}
python scripts/feature_selection/feature_selection.py \
    --config configs/feature_selection/feature_selection_intervention_layer25.yaml
python scripts/intervention/build_feature_set_pool.py \
    results/feature_selection_intervention_layer25/20261105T025730344698Z \
    results/feature_selection_intervention_layer25/20261105T025730344698Z_pool
\end{verbatim}
}
\noindent Final output:
\path{results/feature_selection_intervention_layer25/20261105T025730344698Z_pool/feature_sets.json}.

\paragraph{6. Intervention.} The edit size is chosen on the pilot before the
run is frozen. Each shard runs on its own GPU.

{\footnotesize
\begin{verbatim}
C=configs/intervention/intervention.yaml
python scripts/intervention/run_intervention.py choose-r --config $C
python scripts/intervention/run_intervention.py prepare  --config $C
python scripts/intervention/run_intervention.py run      --config $C --shard <K>
python scripts/intervention/run_intervention.py merge    --config $C
\end{verbatim}
}
\noindent Final output:
\path{results/intervention_layer25/20261106T000808809026Z-02871d677287/merge/}.

\paragraph{7. Analysis.} The edit-size pilot, the main results and the
feature interpretation are in
\path{notebooks/intervention/choose_edit_size_pilot_analysis.ipynb},
\path{notebooks/intervention/full_study_results.ipynb} and
\path{notebooks/intervention/feature_interpretation.ipynb}. The interpretation
reads \path{results/feature_interpretation_layer25/20261106T221912476834Z},
written by \path{scripts/intervention/interpret_features.py}. The exploratory
LLM-edited SAE is analyzed in
\path{notebooks/intervention/llm_edited_sae_analysis.ipynb}.

\section{Building This Report}

{\footnotesize
\begin{verbatim}
latexmk -cd -pdf -interaction=nonstopmode -halt-on-error docs/final_report/main.tex
\end{verbatim}
}
\noindent The PDF is written to \path{docs/final_report/main.pdf}. The figures
are included as PDF files from \path{docs/final_report/figures/} and
\path{docs/final_report/build_figures/}.
\par\endgroup
